\documentclass[french]{article}
\usepackage[utf8]{inputenc}
\usepackage[T1]{fontenc}
\usepackage{lmodern}
\usepackage{geometry}
\usepackage{graphicx}
\usepackage{babel}
\usepackage{amsmath}
\usepackage{amsthm}
\usepackage{amssymb}
\usepackage{hyperref}
\usepackage{stmaryrd}
\usepackage{enumitem}
\usepackage{tcolorbox}
\usepackage{dsfont}
\usepackage{multirow}
\usepackage{etoc}
\usepackage{titlesec}
\usepackage{esint}
\usepackage{cancel}
\usepackage{mathdots}

\setlist[itemize]{label=$\bullet$}


\renewcommand{\P}{\mathbb{P}}
\newcommand{\N}{\mathbb{N}}
\newcommand{\Z}{\mathbb{Z}}
\newcommand{\Q}{\mathbb{Q}}
\newcommand{\R}{\mathbb{R}}
\newcommand{\C}{\mathbb{C}}
\newcommand{\K}{\mathbb{K}}
\renewcommand{\L}{\mathbb{L}}
\newcommand{\U}{\mathbb{U}}
\newcommand{\st}{^*}
\newcommand{\tq}{\middle|}
\newcommand{\inv}{^{-1}}
\newcommand{\E}{\mathbb{E}}
\newcommand{\F}{\mathbb{F}}
\newcommand{\V}{\mathbb{V}}
\renewcommand{\ker}{\text{Ker}}
\newcommand{\im}{\text{Im}}
\newcommand{\card}{\text{Card}}
\newcommand{\Tr}{\text{Tr}}

\title{TD Euclidiens\\ Exercice 3}
\date{}
\begin{document}
\maketitle

\section{Énoncé}
Soit $E$ un espace préhilbertien réel. Pour tout famille finie $(x_i)_{1\leqslant i\leqslant p}$, on pose : $$G(x_1,\ldots,x_p)=((x_i|x_j))_{1\leqslant i,j\leqslant p}\ \ \text{et} \ \ \text{Gram}(x_1,\ldots,x_p)=\det G(x_1,\ldots,x_p)$$
\begin{enumerate}
	\item Montrer que $(x_i)$ est liée si, et seulement si, $\text{Gram}(x_i)=0$.
	\item Montrer qu'ajouter à l'un des $x_i$ une combinaison linéaire des autres ne change pas le rang de $G(x_i)$.
	\item Montrer que $\text{rg}G(x_i)=\text{rg}(x_i)$.
	\item On suppose $E$ muni d'une base $\mathcal{B}=(e_1,\ldots,e_p)$. Donner une relation entre $G(x_i)$, $G(e_i)$ et $M_\mathcal{B}(x_i)$.
	\item Montrer que $G(x_i)$ est positive et en déduire que $\text{Gram}(x_i)\geqslant0$.
	\item Si $F$ est un sous-espace vectoriel dont $y_1$, ..., $y_p$ est une base, on a la formule : $$d(x,F)^2=\frac{\text{Gram}(x,y_1,\ldots,y_p)}{\text{Gram}(y_1,\ldots,y_p)}$$
	\item Lorsque $E$ est de dimension finie $n$, exprimer $\text{Gram}(x_i)$ en fonction du produit mixte des vecteurs $x_i$ pour une orientation quelconque de $E$.
\end{enumerate}
\section{Solution}
Hyper classique ! Peut servir dans d'autres exercices.
\begin{enumerate}
	\item Raisonner en termes de produit scalaire.
	\item Faire des opérations élémentaires.
	\item Utiliser la question précédente.
	\item $G(x_i)=M^TG(e_i)M$ en faisant le calcul des coefficients en exprimant dans la base $\mathcal{B}$.
	\item Calculer comme dans la question précédente.
	\item Commencer par le cas où $x\in F^\bot$ puis utiliser la projection orthogonale qui s'exprime comme $x$ plus une combinaison linéaire des $y_i$.
	\item Le déterminant de Gram est le carré du produit mixte des vecteurs en utilisant la question 4.
\end{enumerate}
\end{document}
